% Part: lambda-calculus
% Chapter: lambda-definability
% Section: lists

\documentclass[../../../include/open-logic-section]{subfiles}

\begin{document}

\olfileid[hi]{lam}{ldf}{lst}
\olsection{सूचियाँ}

सूची $[M_1, M_2, \ldots, M_n]$ को इस प्रकार परिभाषित किया जा सकता है:
\[
  [M_1, M_2, \ldots, M_n] \ident \lambd[sz][s M_1 (s M_2 (\ldots (s M_n z)))]
\]
अनौपचारिक रूप से, किसी सूची को निरूपित करने वाला पद उस सूची का ``संचायक''
फलन है: ऐसा फलन जो दो पद स्वीकार करता है। पहला ऐसा फलन है जो संचित मान के
साथ नया अवयव लेकर नया संचित मान लौटाता है; दूसरा आरंभिक संचित मान है। वह
अवयवों को एक-एक करके संचित करता है और अंत में परिणाम लौटाता है।

\begin{digress}
  यदि आप फलनीय प्रोग्रामिंग से परिचित हैं, तो यह परिभाषा \emph{दाएँ फोल्ड
  (right fold)} जैसी लगेगी: सूची को उसमें समाहित अवयवों के दाएँ फोल्ड फलन
  के रूप में परिभाषित किया गया है।
\end{digress}

इस कूटन के आधार पर हम कुछ उपयोगी फलन सरलता से परिभाषित कर सकते हैं:
\begin{align*}
  \fn{Sum} & \ident
  \lambd[l][l \, (\lambd[xa][\fn{Add} \, x \, a])\,  \num{0}]\\
  \fn{Len} & \ident
  \lambd[l][l \, (\lambd[xa][\fn{Add} \, a \, \num{1}]) \num{0}]
\end{align*}

$\fn{Sum}$ चर्च अंकों की किसी सूची का योग निकालता है। यह सूची पर संचय करता
है, जिसमें आरंभिक मान $\num{0}$ है और प्रत्येक अवयव के लिए नया मान
$\fn{Add} \, x\, a$ लौटाया जाता है (जहाँ $x$ वर्तमान अवयव और $a$ संचित
मान है)। परिणाम सभी अवयवों का योग है।

\begin{prob}
  $\fn{Len}$ क्या करता है? समझाएँ।
\end{prob}

\begin{prob}
  किसी सूची को उलटने वाला फलन परिभाषित करें।
\end{prob}
\end{document}

